\section{Experimental protocol}

\paragraph{Schedules and selection.}
Part~1 sizes begin at 200 and grow geometrically until the exact full cleaned pool. Each dataset-model pair receives ten predefined configurations at 8,000 labels with tuning seed 777. Selection maximizes validation macro-F1, then prefers lower validation inference latency and smaller serialized size. The selected configuration is frozen across final sizes and seeds. Test metrics are absent from tuning records.

The bonus tunes four SupCon and eight DB-CDSupCon settings per target domain. The primary selector is target validation macro-F1, followed by mean source validation macro-F1 and fit time. Pooled CE has no contrastive hyperparameters. The final seeds are 13, 42, and 2026.

\paragraph{Metrics and crossover.}
Macro-F1 is primary. We report mean, sample standard deviation, and seed count for every condition. A Part~1 crossover occurs when the best complete neural mean exceeds the best complete TF-IDF mean by at least 0.01. A sustained crossover satisfies the criterion at that size and every larger tested size. Bonus crossovers apply the same margin against one fixed Part~1 TF-IDF model selected by three-seed mean at each target and size. New target-label savings are reported only when both the target-only and candidate thresholds are observed.

Learning-curve area is the trapezoidal integral over $\log_2$ training size, divided by the tested log-size range. Pareto efficiency uses measured macro-F1 with either fit time or end-to-end latency. Pairwise bonus effects use aligned seed differences and remain descriptive.

\paragraph{Cost measurement.}
The completed \Result{audit_jobs} planned jobs ran on one hardware ID: an NVIDIA RTX 4070 host with an Intel i9-10900K and 31\,GiB RAM. Classical fits used four CPU threads; neural fits used CUDA BF16 autocast with FP32 parameters and optimizer state. Inference ran outside training autocast. Latency uses ten warm-ups and fifty timed calls. Classifier-only latency starts from prepared features; end-to-end latency includes tokenization or encoding. CPU TF-IDF and CUDA neural latency therefore answer a deployment question, not a same-device microbenchmark.

\paragraph{Audit.}
All planned jobs completed, every final cell contains \Result{audit_final_seeds} seeds, and run IDs are unique. The reporting audit found one derived-data exception: the original bonus aggregate selected a TF-IDF winner independently within each seed. The reporting rebuild instead selects one TF-IDF model by the three-seed Part~1 mean, then reuses its aligned seed rows. Atomic measurements are unchanged.
